\chapter{Objections}
\label{ch:essay-objections}

A reasonable objection to the framework developed in this essay is that historical linguistics has long acknowledged that some reconstructions remain uncertain and that not all aspects of linguistic history can be recovered with equal confidence. Comparative linguists distinguish between secure reconstructions supported by regular correspondences across multiple daughter languages and speculative reconstructions based on limited or ambiguous evidence. The field has developed explicit notation for marking reconstruction confidence, with asterisks indicating unattested reconstructed forms and question marks or parentheses indicating uncertain elements. Methodological discussions routinely address the limitations of the comparative method, acknowledging that sound change can obscure original distinctions, that borrowing can introduce forms that resist genealogical explanation, and that chance resemblances can produce spurious cognate sets. Given this existing recognition of reconstruction uncertainty, the objection runs, what does formalization through the admissibility framework add beyond restating in mathematical notation what practitioners already know from experience?

The objection gains force from observing that the practical work of historical reconstruction has proceeded successfully for two centuries without requiring theorems about non-identifiability or diagnostic criteria for distinguishing structural limits from correctable errors. Indo-European comparative linguistics established sound laws, reconstructed proto-forms, and built phylogenetic trees before the formal apparatus developed in this essay existed. The comparative method identifies cognates, projects ancestral states, and assesses reconstruction confidence through procedures that are well understood and widely applied without reference to reachable-set structures or equivalence classes under observation maps. If the field's existing methodological toolkit already handles reconstruction uncertainty effectively, then formalizing that uncertainty as a theorem about observation maps factoring through reachable-set structure might appear to be an exercise in restating the obvious in unnecessarily technical language. The objection deserves serious consideration because it correctly identifies that the general fact of reconstruction uncertainty is not a novel discovery.

The answer to this objection must begin by acknowledging what is correct in it. The theorem does not establish that historical reconstruction is sometimes uncertain; that fact is indeed already known and has been incorporated into linguistic methodology for generations. Theorem~4.2 does not tell historical linguists anything they did not already understand qualitatively about the difficulty of distinguishing sound change from borrowing, the challenges posed by sporadic changes, or the problems created by lexical replacement. The formalization is not intended as a corrective to a field that was previously unaware of these issues. What the theorem provides, and what existing methodology does not, is a principled criterion for diagnosing why a particular reconstruction is uncertain and whether that uncertainty is correctable or fundamental.

The contribution lies specifically in Corollary~4.1, the diagnostic corollary developed in Chapter~4 of \emph{Admissible Reconstruction} and discussed in detail in Chapter~\ref{ch:essay-limits}. The corollary establishes that given a collision between two candidate histories $H_1$ and $H_2$ producing the same or similar observables, one can in principle check whether $\reach{H_1} = \reach{H_2}$. If the reachable sets differ, the reconstruction uncertainty reflects information that the particular observation procedure employed has discarded but that a less lossy procedure could retain. The failure is correctable: gathering different kinds of evidence, applying different comparative techniques, or examining features beyond those currently observed could distinguish the histories. If the reachable sets are equal, the reconstruction uncertainty reflects information that has ceased to exist in any observation preserving only continuation structure. The failure is a structural limit: no amount of additional evidence of the same kind, and no refinement of comparative methodology operating on the same observable features, can resolve the ambiguity.

This distinction is not drawn by existing methodology. When a historical linguist encounters an uncertain reconstruction, current practice might note that the evidence is ambiguous, that multiple hypotheses are consistent with the data, or that the reconstruction should be marked as tentative. What current practice does not provide is a criterion for determining whether the ambiguity is in principle resolvable through additional evidence or whether it reflects a genuine limit of what the observable projection preserves. The difference is not merely academic. If an uncertain reconstruction reflects a correctable observational gap, then the appropriate response is to seek additional evidence: examine more daughter languages, investigate dialectal variation, or incorporate non-phonological data such as morphological alternations or semantic distributions. If it reflects a structural limit, then such efforts are unlikely to succeed, and the appropriate response is to acknowledge the ambiguity as irresolvable rather than treating it as a temporary gap awaiting future discovery.

Consider a concrete example. Suppose a historical linguist reconstructing Proto-Polynesian encounters a lexical item that appears in some daughter languages with initial \textit{p} and in others with initial \textit{f}, and suppose two hypotheses are proposed: either the proto-form had \textit{p} with subsequent lenition to \textit{f} in some branches, or the proto-form had \textit{f} with subsequent fortition to \textit{p} in others. Current methodology would examine the broader pattern of sound correspondences, look for conditioning environments, and assess which hypothesis better fits the overall phonological system, but it might conclude that the evidence is insufficient to decide with certainty. The diagnostic corollary provides an additional step: check whether the two hypothesized histories lead to states with the same reachable-set structure. If both reconstructed proto-forms license the same set of admissible phonological developments in the daughter languages, then the two histories share a reachable set, and Theorem~4.2 applies: no observation map factoring through phonological continuation structure can distinguish them. The uncertainty is structural. If the reachable sets differ—if, for instance, reconstructing \textit{p} predicts different continuation patterns in unstudied daughter languages than reconstructing \textit{f}—then the uncertainty is correctable, and the appropriate response is to examine those unstudied languages or other evidence that would fall within the differential reachable sets.

The objection is therefore correct that the field already knows reconstruction is sometimes uncertain, but it underestimates what the formalization provides. The theorem does not merely restate that uncertainty exists; it provides a diagnostic tool for classifying uncertainty into correctable and structural categories, and it establishes that this classification is in principle decidable by examining reachable-set structure. This transforms reconstruction uncertainty from an undifferentiated acknowledgment that evidence may be insufficient into a structured question with a principled answer: is this case uncertain because we have not yet gathered the right evidence, or because the evidence required to decide has ceased to exist in the observable projection? The formalization thereby adds precision to a distinction that existing methodology gestures toward but does not make explicit or decidable.
